
\documentclass
[   twoside=false,     % Einseitiger oder zweiseitiger Druck?
    fontsize=12pt,     % Bezug: 12-Punkt Schriftgröße
    DIV=15,            % Randaufteilung, siehe Dokumentation "KOMA"-Script
    BCOR=17mm,         % Bindekorrektur: Innen 17mm Platz lassen. Copyshop-getestet.
%    headsepline,
    headsepline,  % Unter Kopfzeile Trennlinie (aus: headnosepline)
    footsepline,  % Über Fußzeile Trennlinie (aus: footnosepline)
    open=right,        % Neue Kapitel im zweiseitigen Druck rechts beginnen lassen
    paper=a4,          % Seitenformat A4
    abstract=true,     % Abstract einbinden
    listof=totoc,      % Div. Verzeichnisse ins Inhaltsverzeichnis aufnehmen
    bibliography=totoc,% Literaturverzeichnis ins Inhaltsverzeichnis aufnehmen
    titlepage,         % Titelseite aktivieren
    headinclude=true,  % Seiten-Head in die Satzspiegelberechnung mit einbeziehen
    footinclude=false, % Seiten-Foot nicht in die Satzspiegelberechnung mit einbeziehen
    numbers=noenddot   % Gliederungsnummern ohne abschließenden Punkt darstellen
]   {scrreprt}         % Dokumentenstil: "Report" aus dem KOMA-Skript-Paket

\usepackage[active]{srcltx}
%\usepackage[activate=normal]{pdfcprot} % Optischer Randausgleich -> pdflatex!
\usepackage{ifthen}
\usepackage[ngerman]{babel}   % Neue Deutsche Rechtschreibung
%\usepackage[latin1]{inputenc} % Zeichencodierung nach ISO-8859-1
\usepackage[utf8]{inputenc}   %	Zeichencodierung nach UTF-8 (Unicode)
\usepackage[T1]{fontenc}
%\usepackage{ae} % obsolet und durch lmodern ersetzt
\usepackage{lmodern}
\usepackage[T1]{url}
\usepackage[final]{graphicx}
%\usepackage[automark]{scrpage2}{scrlayer-scrpage}
\usepackage[automark]{scrlayer-scrpage}
\usepackage{setspace}
%\usepackage[first,light]{draftcopy} % Für Probedruck
\usepackage[plainpages=false,pdfpagelabels,hypertexnames=false]{hyperref}

% Tiefe der Kapitelnummerierung beeinflussen
\setcounter{secnumdepth}{3} % Tiefe der Nummerierung
\setcounter{tocdepth}{3}    % Tiefe des Inhaltsverzeichnisses

% Hier in die zweite geschweifte Klammer jeweils
% die persönlichen Daten und das Thema der Arbeit eintragen:
\newcommand{\artderausarbeitung}{Hausarbeit}
\newcommand{\namedesautors}{Jonathan Skopp}
\newcommand{\themaderarbeit}{Selbstreflexion zum Modul Soft Skills}

% PDF Metadaten definieren
\hypersetup{
   pdftitle={\themaderarbeit},
   pdfsubject={\artderausarbeitung},
   pdfauthor={\namedesautors},
   pdfkeywords={\artderausarbeitung; TU-Ilmenau; ZIB;}}


% Abkürzungsverzeichnis beeinflussen. Hier nichts ändern!
\usepackage[intoc]{nomencl}
  \AtBeginDocument{\setlength{\nomlabelwidth}{.25\columnwidth}}
  \let\abbrev\nomenclature
  \renewcommand{\nomname}{Abkürzungsverzeichnis und Formelzeichen}
  \renewcommand{\nomlabel}[1]{#1 \dotfill}
  \setlength{\nomitemsep}{-\parsep}
  \makenomenclature

\usepackage[normalem]{ulem}
  \newcommand{\markup}[1]{\textbf{#1}}

% Seitenlayout festlegen. Hier nichts ändern!
\pagestyle{scrplain}
\ihead[]{\headmark}
\ohead[]{\pagemark}
\chead[]{}
\ifoot[]{}
\ofoot[]{\scriptsize \artderausarbeitung\ \namedesautors}
\cfoot[]{}
\renewcommand{\titlepagestyle}{scrheadings}
\renewcommand{\partpagestyle}{scrheadings}
\renewcommand{\chapterpagestyle}{scrheadings}
\renewcommand{\indexpagestyle}{scrheadings}

% Abschnittsweise Nummerierung anstatt fortlaufend. Hier nichts ändern!
\makeatletter
\@addtoreset{equation}{chapter}
\@addtoreset{figure}{chapter}
\@addtoreset{table}{chapter}
\renewcommand\theequation{\thechapter.\@arabic\c@equation}
\renewcommand\thefigure{\thechapter.\@arabic\c@figure}
\renewcommand\thetable{\thechapter.\@arabic\c@table}
\makeatother

% Quelltextrahmen, klein. Hier nichts ändern!
\newsavebox{\inhaltkl}
\def\rahmenkl{\sbox{\inhaltkl}\bgroup\small\renewcommand{\baselinestretch}{1}\vbox\bgroup\hsize\textwidth}
\def\endrahmenkl{\par\vskip-\lastskip\egroup\egroup\fboxsep3mm%
\framebox[\textwidth][l]{\usebox{\inhaltkl}}}

% Quelltextrahmen, normale Groesse. Hier nichts ändern!
\newsavebox{\inhalt}
\def\rahmen{\sbox{\inhalt}\bgroup\renewcommand{\baselinestretch}{1}\vbox\bgroup\hsize\textwidth}
\def\endrahmen{\par\vskip-\lastskip\egroup\egroup\fboxsep3mm%
\framebox[\textwidth][l]{\usebox{\inhalt}}}

% Trennvorschläge für falsch getrennte Wörter.
% Wird häufig bei eingedeutschen Wörtern benötigt, da LaTeX hierbei
% gerne falsch trennt. Alternativ kann auch im Fliesstext ein
% Trennvorschlag per "\-" hinterlegt werden, bspw.:
% Die Hard\-ware besteht aus A und B.
\hyphenation{
Hard-ware
}

% Sonstige Befehlsdefinitionen hier ablegen.
\newcommand{\entspricht}{\stackrel{\wedge}{=}}

% Tabellenspaltendefinitionen mit fester Breite --> somit Zeilenumbruch innerhalb einer Zelle möglich
% aus http://www.torsten-schuetze.de/tex/tabsatz-2004.pdf
\usepackage{array, booktabs}
\newcolumntype{f}{>{$}l<{$}}
\newcolumntype{n}{>{\raggedright}l}
\newcolumntype{N}{>{\scriptsize}l}
\newcolumntype{v}[1]{>{\raggedright\hspace{0pt}}m{#1}}
\newcolumntype{V}[1]{>{\scriptsize\raggedright\hspace{0pt}}m{#1}}
\newcolumntype{Z}[1]{>{\raggedright\centering}m{#1}}
\newcolumntype{k}[1]{>{\raggedright}p{#1}}
% ergibt Tabllenspalte fester Breite, linksbündig
% Umbruch innerhalb der Zelle mit \\, neue Tabellezeile mit \tabularnewline
% \addlinespace für Gruppentrennung (aus \texttt{booktabs.sty})


\begin{document}
\onehalfspacing
%% ++++++++++++++++++++++++++++++++++++++++++++++++++++++++++++
%% Hauptdatei, Wurzel des Dokuments
%% ++++++++++++++++++++++++++++++++++++++++++++++++++++++++++++
%
%  Gerüst:
%  * Version 0.13
%  * M. Sc. Silvia Krug, silvia.krug@tu-ilmenau.de
%  * Fachgebiet Kommunikationsnetze, TU Ilmenau
%
%  Für Hauptseminare, Studienarbeiten, Diplomarbeiten
%
%  Autor           : Max Mustermann
%  Letzte Änderung : 04.12.2015
%

\begin{titlepage}
\centering
\includegraphics[scale=0.5]{bilder/tui_logo}\\[3ex]
{\Large \textsc{Technische Universität Ilmenau}}\\[3ex]
{\Large Fakultät für Elektrotechnik und Informationstechnik}\\[3ex]
\vfill
{\Large \textbf{\artderausarbeitung}}\\[4ex]
{\large \textbf{\themaderarbeit}}\\[5ex]
\vfill
\begin{tabular}{rl}
\hline\\
vorgelegt von:          & \quad \namedesautors\\[1,5ex]
eingereicht am:         & \quad 31.\,12.\,2011\\[1,5ex]
geboren am:             & \quad 31.\,12.\,1985 in Ilmenau\\[1,5ex]
Studiengang:            & \quad Ingenieurinformatik\\[1,5ex]
Studienrichtung:        & \quad Multimediale Informations- und\\[1,5ex]
                        & \quad Kommunikationssysteme\\[5ex]
Anfertigung im Fachgebiet:
                        & \quad Kommunikationsnetze\\[1,5ex]
                        & \quad Fakultät für Elektrotechnik und Informationstechnik\\[1,5ex]
Verantwortlicher Professor:
                        & \quad Prof.~Dr.~rer.~nat.~Jochen Seitz\\[1,5ex]
Wissenschaftlicher Betreuer:
                        & \quad M.~Sc.~Vorname Nachname
\end{tabular}
\vfill
\end{titlepage}










% Inhaltsverzeichnis
\cleardoublepage % Seitenumbruch erzwingen vor Änderung des Nummerierungsstils
\pagenumbering{roman} % Nummerierung der Seiten ab hier: i, ii, iii, iv...
\pagestyle{scrheadings} % Ab hier mit Kopf- und Fusszeile
\tableofcontents

% Die einzelnen Kapitel
\cleardoublepage % Seitenumbruch erzwingen vor Änderung des Nummerierungsstils
\pagenumbering{arabic} % Nummerierung der Seiten ab hier: 1, 2, 3, 4...
\chapter{Was habe ich mitgenommen ? }

Interessanter Weise habe ich viel informatives Mitgenommen. Durch die Vorlesungen konnte man sehr viel theoretisches Wissen aufnehmen. Durch die praktische Anwendung konnte ich dieses auch sehr gut Umsetzen. Alle Teilnehmer waren also "gezwungen" die gestellte Aufgabe zu lösen. Obwohl es sehr anstrengend war, hat es viel Spaß gemacht. 




\section{Rückblick}
\subsection{Gute Aspekte}
Mir hat das Modul Spaß gemacht. Das liegt vor allem daran, dass es mal etwas anderes war. Die schnelle Digitalisierung die Corona erfordert hat, wurde sehr gut umgesetzt. 


\subsection{Schlechte Aspekte}
Ich fande nichts in diesem Modul als schlecht. Ich würde die Verlinkungen etwas anpassen. So muss man sich in Moodle anmelden um die Prezi zu suchen um dann dort die Links zu den Videos zu finden. Eine Idee wäre es, diese über verschiedene Chapter im Moodle direkt zu publizieren. Es wäre auch schön gewesen, wenn man die Prezi hätte herunterladen können.

\chapter{Was wusste ich bereits}

\begin{quote}
	Test
\end{quote}
\chapter{Neues Wissen}

Hallo Welt
\chapter{Resonanz}

Hallo Welt


%\appendix
%\part*{Anhang}


% Abschlusserklärung

\chapter*{Erklärung}
\addcontentsline{toc}{chapter}{Erklärung}
\ihead[]{Erklärung}

Die vorliegende Arbeit habe ich selbstständig ohne Benutzung anderer als der
angegebenen Quellen angefertigt. Alle Stellen, die wörtlich oder sinngemäß
aus veröffentlichten Quellen entnommen wurden, sind als solche
kenntlich gemacht. Die Arbeit ist in gleicher oder ähnlicher Form oder
auszugsweise im Rahmen einer oder anderer Prüfungen noch nicht vorgelegt
worden.
\\[2cm]
Ilmenau, den 15.\,07.\,2020\hfill \namedesautors


\end{document}

